\documentclass[11pt]{article}
\usepackage[T1]{fontenc}
\usepackage{amsmath}
\usepackage{booktabs}
\usepackage{xcolor}
\input{comparison.tex}

\title{Three-Hour Generator Scheduling}
\author{Synthetic training example}
\date{}

\begin{document}
\maketitle

\begin{abstract}
We enumerate feasible integer generator outputs for three independent hours.
The base schedule costs \BaseCost{} cost units. Raising demand by 10 MW in
each hour increases the cost to \ScenarioCost{}. A demand of 141 MW in hour 2
exceeds available capacity.
\end{abstract}

\section{Data and model}
The training case provides generator limits and prices in \texttt{units.csv}
and hourly demand in \texttt{demand.csv}. For each unit \(i\) and hour \(t\),
output \(p_{i,t}\) is either zero or an integer in the inclusive range
\([\ell_i,u_i]\) MW. The cost is
\(f_i\mathbf{1}_{\{p_{i,t}>0\}}+c_i p_{i,t}\).
We minimize the sum of both units' costs subject to
\(p_{\mathrm{A},t}+p_{\mathrm{B},t}=d_t\). With independent hours, we enumerate
unit A's allowed outputs, compute B's residual output, check the bounds and
select the cheapest feasible pair.

\section{Base schedule}
Table~\ref{tab:base} gives the result for Question 1. The total cost is
\BaseCost{} cost units.

\begin{table}[ht]
\centering
\caption{Minimum-cost schedule under the original demand.}
\label{tab:base}
\input{base-schedule.tex}
\end{table}

\section{Demand scenario}
For Question 2, Table~\ref{tab:scenario} gives the updated schedule.
The total cost increases by \CostIncrease{} cost units
(Figure~\ref{fig:costs}).

\begin{table}[ht]
\centering
\caption{Minimum-cost schedule with demand raised by 10 MW in each hour.}
\label{tab:scenario}
\input{scenario-schedule.tex}
\end{table}

\begin{figure}[ht]
\centering
\begin{tabular}{lp{48mm}r}
Base & {\color{blue!65!black}\rule{\BaseBar}{2.5mm}} & \BaseCost \\
+10 MW per hour & {\color{teal!65!black}\rule{\ScenarioBar}{2.5mm}} & \ScenarioCost
\end{tabular}
\caption{Total cost in both scenarios, starting at zero (cost units).}
\label{fig:costs}
\end{figure}

\section{Independent checks and scope}
When demand is 95 MW, both units must operate. A costs less per MW, so its
output is 80 MW and B supplies the remaining 15 MW. The hourly cost is
\(60+4(80)+30+6(15)=500\) cost units. A putative 85 MW from A and 10 MW from
B meets demand but violates A's upper limit.

For Question 3, the two maximum outputs sum to \(80+60=140\) MW.
The 141 MW demand is therefore infeasible, and the script rejects that input.
The results apply to the stated independent hours and integer MW resolution.
Inter-hour ramping and start-up constraints require a different model.

\end{document}
